\documentclass[a4paper,11pt]{article}
\usepackage{fontspec}
\usepackage{unicode-math}
\usepackage{amsmath,amssymb,amsthm}
\usepackage{longtable,booktabs,array}
\usepackage{hyperref}
\usepackage[margin=25mm]{geometry}
\setmainfont{Latin Modern Roman}

\newtheorem{theorem}{Theorem}[section]
\newtheorem{proposition}[theorem]{Proposition}
\newtheorem{lemma}[theorem]{Lemma}
\newtheorem{corollary}[theorem]{Corollary}
\newtheorem{definition}[theorem]{Definition}
\newtheorem{remark}[theorem]{Remark}
\newtheorem{observation}[theorem]{Observation}

\providecommand{\tightlist}{\setlength{\itemsep}{0pt}\setlength{\parskip}{0pt}}
\newcounter{none}


\begin{document}


\section{D1 v2.0: Particle Data Structure Specification of the Central
Projection Framework --- Frequency Interpretation and Dynamics
Extension}\label{d1-v2.0-particle-data-structure-specification-of-the-central-projection-framework-frequency-interpretation-and-dynamics-extension}

\textbf{Author}: Noriaki Kihara \textbf{Affiliation}: WF System Co.,
Ltd. \textbf{ORCID}: 0009-0004-6753-4020 \textbf{Version}: v2.1
\textbf{Date}: May 2026 \textbf{DOI}:
\href{https://doi.org/10.5281/zenodo.19941250}{10.5281/zenodo.19941250}
(v2.0; will be updated when v2.1 is published) \textbf{Concept DOI}:
\href{https://doi.org/10.5281/zenodo.19941249}{10.5281/zenodo.19941249}
\textbf{Predecessor}: v1.0 (April 2026, unpublished) / v2.0 (May 2026
published)

\begin{center}\rule{0.5\linewidth}{0.5pt}\end{center}

\subsection{§0 Position of This Paper}\label{position-of-this-paper}

This paper is a design specification, not a mathematical proof. It
explicitly defines, as a typed-integer-arithmetic specification, the
particle data structures that have been implicitly assumed in the
preceding paper series {[}W1{]}--{[}W11{]}, {[}M1{]}--{[}M3{]},
supplementary lectures {[}Sup1{]}--{[}Sup4{]}, and the thought
experiment {[}TE\_Raxis{]}.

\subsubsection{Revisions from v2.0
(v2.1)}\label{revisions-from-v2.0-v2.1}

\textbf{Removal of \texttt{δδt} field}: The \texttt{δδt\ :\ DeltaT}
field added to the Particle structure in v2.0 (processing-period change,
acceleration equivalent) is \textbf{removed} in v2.1. Reasons:

\begin{enumerate}
\def\labelenumi{\arabic{enumi}.}
\tightlist
\item
  \textbf{Reversibility violation}: A \texttt{step(p)} containing
  \texttt{δt\ +=\ δδt} does not return to its original state when δt's
  sign is reversed (violates the information conservation requirement of
  {[}W3{]}).
\item
  \textbf{Inconsistency with conservation laws}: Changing δt for a
  single particle without external interaction means the proper time
  changes ex nihilo, implicitly breaking energy and momentum
  conservation.
\item
  \textbf{GR free-particle principle}: A free particle follows a
  geodesic, and acceleration effects must arise from external forces
  (interactions). The basic operation for a single particle should not
  handle acceleration.
\end{enumerate}

All acceleration effects are represented \textbf{within interaction
processing (between particles or between particle and Mesh)} as changes
in σ or δt that satisfy conservation laws. See the separate paper D2
(``Specification of Dynamics Rules of the Central Projection
Framework,'' in preparation).

The corrected Particle structure has \textbf{4 fields}
\texttt{(σ,\ V,\ δV,\ δt)} (§5). The corresponding \texttt{δδt\_mesh} is
also removed from the Mesh structure (§7).

\subsubsection{Major Revisions from v1.0 (introduced in v2.0, continued
in
v2.1)}\label{major-revisions-from-v1.0-introduced-in-v2.0-continued-in-v2.1}

\begin{enumerate}
\def\labelenumi{\arabic{enumi}.}
\tightlist
\item
  \textbf{Reinterpret σ.k\_xyzt as ``vibrational frequencies of each
  axis''} rather than ``wave numbers / acceleration'' (§1.4)
\item
  \textbf{Add δV (phase displacement) to the Particle structure} (§5).
  \textasciitilde v2.0 also added δδt, but this was removed in
  v2.1.\textasciitilde{}
\item
  \textbf{Introduce the Mesh structure to separate longitudinal modes
  (Higgs, gravitational waves) from Particles} (§7)
\item
  \textbf{Eliminate special status of t-axis, fully preserving 6-axis
  symmetry} (§3, §6)
\end{enumerate}

The data type system (int256-based) and integer-arithmetic principles
(no division, no square root, no trigonometric functions) remain
unchanged from v1.0.

\subsubsection{Claims of This Paper}\label{claims-of-this-paper}

\begin{enumerate}
\def\labelenumi{\arabic{enumi}.}
\tightlist
\item
  The boundedness of the universe and the discreteness at the Planck
  scale together permit all phase quantities to be represented as
  256-bit signed integers (§1).
\item
  The \([L^{-1}]\) dimension of the xyzt axes is naturally interpreted
  as ``vibrational frequency along each axis in phase space,'' fully
  preserving 6-axis symmetry (§3).
\item
  The complete state of a particle can be described by 4 fields:
  \texttt{SignVector} (identity label) + \texttt{OneDWave{[}κ{]}} (phase
  state) + \texttt{δV} (phase displacement) + \texttt{δt} (processing
  period) (§5).
\item
  Longitudinal modes such as the Higgs should be represented not as
  Particles but as vibrations of the \texttt{Mesh}, allowing unified
  treatment of mass generation, gravitational waves, and cosmological
  vibrations (§7).
\item
  All operations close under integer addition and multiplication,
  requiring no division, square root, or trigonometric functions (§6).
\item
  The basic operation \texttt{step(p)} (which only updates V) satisfies
  all five operational principles: time-reversal symmetry,
  conservation-law consistency, null transparency, type integrity, and
  INPUT/OUTPUT exchangeability (§6).
\end{enumerate}

\subsubsection{Out of Scope}\label{out-of-scope}

\begin{itemize}
\tightlist
\item
  Dynamical formulation of operation rules (physical meaning of phase
  evolution, concrete forms of interactions) --- deferred to a separate
  paper (D2).
\item
  Numerical derivation of specific physical quantities.
\item
  Hardware implementation details.
\item
  Mathematical proofs of theorems.
\end{itemize}

\subsubsection{Referenced Papers}\label{referenced-papers}

\begin{itemize}
\tightlist
\item
  {[}W3{]} Structural Requirements for a Theory of Everything (DOI:
  10.5281/zenodo.19601592)
\item
  {[}W8{]} Combinatorial Structure of the 6-Dimensional Hyperrectangle
  (DOI: 10.5281/zenodo.19762134)
\item
  {[}W9{]} sine-Gordon Equation --- Foundation of Topological Solitons
  (DOI: 10.5281/zenodo.19650966)
\item
  {[}W10{]} Four Modes of Shape-Invariant Waves (DOI:
  10.5281/zenodo.19709798)
\item
  {[}W11{]} Interactions of Shape-Invariant Waves (DOI:
  10.5281/zenodo.19763463)
\item
  {[}M1{]} Three-Layer Model of Central Projection: R--R₁--R₀ (DOI:
  10.5281/zenodo.19691713)
\item
  {[}TE\_Raxis{]} Reexamination of the 6-Dimensional Sign Vector xyztRQ
  (DOI: 10.5281/zenodo.19904714)
\end{itemize}

\begin{center}\rule{0.5\linewidth}{0.5pt}\end{center}

\subsection{§1 int256 Boundedness and Dimensional
Interpretation}\label{int256-boundedness-and-dimensional-interpretation}

\subsubsection{§1.1 Basis for Boundedness Requirement (unchanged from
v1.0)}\label{basis-for-boundedness-requirement-unchanged-from-v1.0}

{[}W3{]} requires boundedness as a structural condition for any theory
of everything. In the three-layer structure \(R\)--\(R_1\)--\(R_0\) of
{[}M1{]}, the maximum curvature radius \(R_0\) is finite. Combined with
discreteness at the Planck scale, this yields an upper bound on the bits
needed to represent all phase quantities.

\subsubsection{§1.2 Required Number of
Bits}\label{required-number-of-bits}

\textbf{Spatial axis}: From the observable universe diameter
\(D \approx 1.892 \times 10^{27}\) m and Planck length
\(l_P = 1.616 \times 10^{-35}\) m,

\[n_{\text{space}} = D / l_P \approx 1.171 \times 10^{62}, \quad b_{\text{space}} = 207 \text{ bits}\]

\textbf{Time axis}: From cosmic time \(T \approx 8.710 \times 10^{17}\)
s and Planck time \(t_P = 5.391 \times 10^{-44}\) s,

\[n_{\text{time}} = T / t_P \approx 1.616 \times 10^{61}, \quad b_{\text{time}} = 204 \text{ bits}\]

\textbf{Conclusion}: 256-bit signed integers leave 14-digit margin in
space and 15-digit margin in time.

\subsubsection{§1.3 Natural Units and
Dimensions}\label{natural-units-and-dimensions}

Under \(c = 1\), \(G = 1\) (geometrized units; \(\hbar\) is not
adopted),

\[[L] = [T] = [M], \quad [L^{-1}] = [M^{-1}]\]

\subsubsection{§1.4 Frequency Interpretation (Core Revision in
v2.0)}\label{frequency-interpretation-core-revision-in-v2.0}

In v1.0, the \([L^{-1}]\) dimension of the xyzt axes was interpreted
ambiguously as ``wave number'' or ``acceleration.'' In v2.0, this is
unified as \textbf{``vibrational frequency along each axis''} as the
natural interpretation in phase space.

{\def\LTcaptype{none} % do not increment counter
\begin{longtable}[]{@{}
  >{\raggedright\arraybackslash}p{(\linewidth - 6\tabcolsep) * \real{0.2500}}
  >{\raggedright\arraybackslash}p{(\linewidth - 6\tabcolsep) * \real{0.2500}}
  >{\raggedright\arraybackslash}p{(\linewidth - 6\tabcolsep) * \real{0.2500}}
  >{\raggedright\arraybackslash}p{(\linewidth - 6\tabcolsep) * \real{0.2500}}@{}}
\toprule\noalign{}
\begin{minipage}[b]{\linewidth}\raggedright
Quantity
\end{minipage} & \begin{minipage}[b]{\linewidth}\raggedright
Dimension
\end{minipage} & \begin{minipage}[b]{\linewidth}\raggedright
v1.0 Interpretation
\end{minipage} & \begin{minipage}[b]{\linewidth}\raggedright
\textbf{v2.1 Interpretation}
\end{minipage} \\
\midrule\noalign{}
\endhead
\bottomrule\noalign{}
\endlastfoot
σ.k\_x, k\_y, k\_z & \([L^{-1}]\) & spatial wave number / acceleration &
\textbf{vibrational frequency of each spatial axis} \\
σ.k\_t & \([L^{-1}]\) & temporal frequency & \textbf{vibrational
frequency of t-axis} (equal to others) \\
σ.R & \([L]=[M]\) & curvature radius & \textbf{scale quantity /
corresponding period} (dual) \\
σ.Q & \([1]\) & internal symmetry & internal symmetry (unchanged) \\
δt & \([L]=[M]\) & proper time step & \textbf{processing step period}
(not ``time'') \\
\textsubscript{δδt} & \textsubscript{\([L]=[M]\)} & \textasciitilde(not
in v1.0)\textasciitilde{} & \textasciitilde added in v2.0, removed in
v2.1 (violates conservation/reversibility)\textasciitilde{} \\
\end{longtable}
}

\subsubsection{§1.5 Advantages of the Frequency
Interpretation}\label{advantages-of-the-frequency-interpretation}

\textbf{(1) Complete preservation of 6-axis symmetry}: All four
position-type axes (xyzt) carry independent vibrational modes; the
t-axis is on equal footing with the others.

\textbf{(2) Elimination of continuum-mechanical concepts}:
``Acceleration'' is a continuum concept (second time derivative of
position) that does not fit naturally into discrete integer arithmetic.
``Vibrational frequency'' is the natural discrete quantity in phase
space.

\textbf{(3) Acceleration handled as interaction}: Acceleration-like
effects (gravitational fields, inertial frame influences) are not
represented for a single particle but \textbf{entirely as interaction
processing} (between particles or between particle and Mesh), where σ or
δt change in a conservation-consistent manner. This is consistent with
the GR free-particle principle (free particles follow geodesics). See
the separate paper D2.

\textbf{(4) Frequency × period = phase}: \(\sigma.k \cdot \delta t\)
becomes a dimensionless phase increment (\([L^{-1}] \cdot [L] = [1]\)),
making the phase evolution rule type-consistent.

\textbf{(5) Automatic inclusion of special relativity}: The dispersion
relation
\(\sigma.k_t^2 = \sigma.k_x^2 + \sigma.k_y^2 + \sigma.k_z^2 + m^2\) can
be naturally written as integer relations among frequencies.

\begin{center}\rule{0.5\linewidth}{0.5pt}\end{center}

\subsection{§2 Definition of Meta Types}\label{definition-of-meta-types}

\subsubsection{§2.1 Scalar Types (unchanged from
v1.0)}\label{scalar-types-unchanged-from-v1.0}

\begin{verbatim}
Type         Dimension  Range                     Use
L256         [L]=[M]    |v| ≤ 2^255 - 1          length / time / mass / period
Linv256      [L⁻¹]      |v| ≤ 2^255 - 1          frequency / inverse scale
Scalar256    [1]        |v| ≤ 2^255 - 1          dimensionless
FullPhase    [1]        0 ≤ v ≤ P                0–720°
CyclePhase   [1]        0 ≤ v ≤ P/2              0–360° (position phase)
DeltaPhase   [1]        -P/2 ≤ v ≤ P/2           ±360° (phase displacement)
SpinPhase    [1]        v ∈ {0, P/4, P/2, P}     {0, 180, 360, 720}°
\end{verbatim}

Phase constant \(P = \lfloor (2^{255}-1)/720 \rfloor \times 720\).

\subsubsection{§2.2 Axis Component Types (renamed and reinterpreted in
v2.0)}\label{axis-component-types-renamed-and-reinterpreted-in-v2.0}

{\def\LTcaptype{none} % do not increment counter
\begin{longtable}[]{@{}
  >{\raggedright\arraybackslash}p{(\linewidth - 6\tabcolsep) * \real{0.2500}}
  >{\raggedright\arraybackslash}p{(\linewidth - 6\tabcolsep) * \real{0.2500}}
  >{\raggedright\arraybackslash}p{(\linewidth - 6\tabcolsep) * \real{0.2500}}
  >{\raggedright\arraybackslash}p{(\linewidth - 6\tabcolsep) * \real{0.2500}}@{}}
\toprule\noalign{}
\begin{minipage}[b]{\linewidth}\raggedright
Type
\end{minipage} & \begin{minipage}[b]{\linewidth}\raggedright
Base
\end{minipage} & \begin{minipage}[b]{\linewidth}\raggedright
Dimension
\end{minipage} & \begin{minipage}[b]{\linewidth}\raggedright
Physical Meaning (v2.0)
\end{minipage} \\
\midrule\noalign{}
\endhead
\bottomrule\noalign{}
\endlastfoot
\texttt{SpatialFreq} \textless: Linv256 & int256 & \([L^{-1}]\) &
vibrational frequency of x, y, z axes \\
\texttt{TemporalFreq} \textless: Linv256 & int256 & \([L^{-1}]\) &
vibrational frequency of t-axis (equal to others) \\
\texttt{CurvatureR} \textless: L256 & int256 & \([L]\) & scale / period
(signed square number) \\
\texttt{ColorQ} \textless: Scalar256 & int256 & \([1]\) & internal
symmetry label (finite group) \\
\texttt{DeltaT} \textless: L256 & int256 & \([L]\) & processing step
period \\
\end{longtable}
}

\textbf{Renaming from v1.0}: \texttt{SpatialK} → \texttt{SpatialFreq},
\texttt{TemporalK} → \texttt{TemporalFreq}. The type system itself is
unchanged; only the interpretation is clarified.

\subsubsection{§2.3 Operation Rules (unchanged from
v1.0)}\label{operation-rules-unchanged-from-v1.0}

\textbf{Addition}: Only between same types (CyclePhase + DeltaPhase =
CyclePhase, etc.).

\textbf{Multiplication}: - \texttt{L256\ ×\ L256\ →\ L2\_256} \([L^2]\)
- \texttt{Linv256\ ×\ Linv256\ →\ Linv2\_256} \([L^{-2}]\) -
\texttt{Linv256\ ×\ L256\ →\ Scalar256} \([1]\) (dimensionless)

\textbf{Division}: Not used.

\begin{center}\rule{0.5\linewidth}{0.5pt}\end{center}

\subsection{§3 Definition of SignVector (Frequency Interpretation
Version)}\label{definition-of-signvector-frequency-interpretation-version}

\subsubsection{§3.1 Structure}\label{structure}

\begin{verbatim}
SignVector = (
    x : SpatialFreq,    -- frequency of x-axis vibration [L⁻¹]
    y : SpatialFreq,    -- frequency of y-axis vibration
    z : SpatialFreq,    -- frequency of z-axis vibration
    t : TemporalFreq,   -- frequency of t-axis vibration [L⁻¹] (equal to others)
    R : CurvatureR,     -- scale quantity / period [L]
    Q : ColorQ,         -- internal symmetry label [1]
)
\end{verbatim}

\subsubsection{§3.2 Physical Meaning}\label{physical-meaning}

Each σ component represents either a ``vibrational frequency along the
corresponding axis'' or a ``scale quantity'':

\begin{itemize}
\tightlist
\item
  σ.k\_axis = 0: the axis does not vibrate (no structure in that
  direction)
\item
  σ.k\_axis ≠ 0: the axis vibrates at frequency
  \textbar k\_axis\textbar{}
\item
  sign(k\_axis): direction of phase rotation (+ for clockwise, − for
  counterclockwise)
\item
  \textbar k\_axis\textbar: magnitude of vibrational frequency (discrete
  integer value)
\end{itemize}

\subsubsection{§3.3 Invariants}\label{invariants}

The values of SignVector change in interaction events, but specific
combinations define stable modes:

\begin{itemize}
\tightlist
\item
  \(\kappa_s = |\{i \in \{x,y,z\} : k_i \neq 0\}|\) --- number of
  non-zero spatial axes
\item
  \(\kappa_t = (1 \text{ if } k_t \neq 0 \text{ else } 0)\) --- non-zero
  indicator for t-axis
\item
  \(\kappa = \kappa_s + \kappa_t \in \{0, 1, 2\}\) --- {[}W10{]}
  Proposition 2.1: \(\kappa \geq 3\) is unstable
\end{itemize}

\subsubsection{§3.4 Dispersion Relation}\label{dispersion-relation}

The mass-shell condition is directly expressed as an integer relation
among frequencies:

\[\sigma.k_t^2 = \sigma.k_x^2 + \sigma.k_y^2 + \sigma.k_z^2 + m^2\]

\begin{itemize}
\tightlist
\item
  \(m = 0\) (photon, gluon): \(k_t^2 = k_x^2 + k_y^2 + k_z^2\)
\item
  \(m \neq 0\) (massive particle): the square root of the difference
  corresponds to mass
\end{itemize}

This automatically incorporates special relativity into the framework.

\subsubsection{§3.5 Correspondence with {[}W8{]} Table A (continued from
v1.0)}\label{correspondence-with-w8-table-a-continued-from-v1.0}

The sign patterns \(\{+, -, 0\}^6\) in {[}W8{]} Table A correspond to
the signs of SignVector. Integer values (axis scale values) correspond
to \(v_i\) in {[}Mass Structure{]} §1.

Representative particles:

{\def\LTcaptype{none} % do not increment counter
\begin{longtable}[]{@{}
  >{\raggedright\arraybackslash}p{(\linewidth - 6\tabcolsep) * \real{0.2500}}
  >{\raggedright\arraybackslash}p{(\linewidth - 6\tabcolsep) * \real{0.2500}}
  >{\raggedright\arraybackslash}p{(\linewidth - 6\tabcolsep) * \real{0.2500}}
  >{\raggedright\arraybackslash}p{(\linewidth - 6\tabcolsep) * \real{0.2500}}@{}}
\toprule\noalign{}
\begin{minipage}[b]{\linewidth}\raggedright
Particle
\end{minipage} & \begin{minipage}[b]{\linewidth}\raggedright
Set \((x,y,z,t,R,Q)\)
\end{minipage} & \begin{minipage}[b]{\linewidth}\raggedright
\(\kappa\)
\end{minipage} & \begin{minipage}[b]{\linewidth}\raggedright
Classification
\end{minipage} \\
\midrule\noalign{}
\endhead
\bottomrule\noalign{}
\endlastfoot
Photon \(\gamma\) & \((0,0,0,0,+,0)\) & 0 & gauge boson (free, spherical
wave) \\
Gluon \(g\) & \((0,0,0,0,0,Q_i \to Q_j)\) & 0 & gauge boson \\
Higgs \(H\) & \((0,0,+,0,0,0)\) & 1 & scalar (longitudinal
interpretation in §7) \\
\(W^+\) & \((+,0,0,0,+,0)\) & 1 & gauge boson \\
\(Z^0\) & \((0,+,0,0,+,0)\) & 1 & gauge boson \\
Graviton & \((0,0,0,\pm,\pm,0)\) & 1 & tensor boson \\
Electron \(e^-\) & \((+,0,0,+,0,0)\) & 2 & fermion \\
\end{longtable}
}

\begin{center}\rule{0.5\linewidth}{0.5pt}\end{center}

\subsection{§4 Definition of OneDWave}\label{definition-of-onedwave}

\subsubsection{§4.1 Structure (unchanged from
v1.0)}\label{structure-unchanged-from-v1.0}

\begin{verbatim}
OneDWave = (
    θ₀ : CyclePhase,   -- position phase [1]   (S1 §2.1, W9 §8.2)
    A  : Scalar256,     -- amplitude [1]        (S1 §2.1, W9 §8.2)
    ℓ  : DeltaPhase,    -- spread phase [1]     (W9 §8.3)
)
\end{verbatim}

\subsubsection{§4.2 Physical Meaning}\label{physical-meaning-1}

Each particle has one OneDWave per non-zero position-type axis (at most
\(\kappa\)). Each OneDWave holds the current phase state of vibration in
that axis direction.

\begin{itemize}
\tightlist
\item
  \(\theta_0\): current phase of vibration (free parameter, evolves over
  time)
\item
  \(A\): amplitude (free parameter)
\item
  \(\ell\): spatial extent of localization (dependent on SignVector,
  {[}W9{]} §8.3)
\end{itemize}

\begin{center}\rule{0.5\linewidth}{0.5pt}\end{center}

\subsection{§5 Definition of Particle (v2.1: 4-field
structure)}\label{definition-of-particle-v2.1-4-field-structure}

\subsubsection{§5.1 Structure}\label{structure-1}

\begin{verbatim}
Particle = (
    σ   : SignVector,    -- identity (vibrational frequencies of each axis)
    V   : OneDWave[κ],   -- current phase state of each non-zero position-type axis
    δV  : OneDWave[κ],   -- phase displacement per processing step (added in v2.0)
    δt  : DeltaT,         -- current processing step period
)
\end{verbatim}

\textbf{Change from v2.0}: The \texttt{δδt\ :\ DeltaT} field is removed
(5 fields → 4 fields). See §0 for reasons.

\subsubsection{§5.2 Physical Meaning of Each
Field}\label{physical-meaning-of-each-field}

\paragraph{\texorpdfstring{\texttt{σ\ :\ SignVector}}{σ : SignVector}}\label{ux3c3-signvector}

The particle's identity. Specifies frequencies of vibration in each
axis. Values change in interaction events (formulated as annihilation +
creation, but as a data structure σ is updated).

\paragraph{\texorpdfstring{\texttt{V\ :\ OneDWave{[}κ{]}}}{V : OneDWave{[}κ{]}}}\label{v-onedwaveux3ba}

Current phase state of each non-zero position-type axis. Array length
\(\kappa\) equals the number of non-zero position-type axes in
SignVector (\(\leq 2\)). Each V{[}j{]} holds the (position phase,
amplitude, spread) of vibration in the corresponding axis.

Index correspondence:

\begin{verbatim}
σ.k_x ≠ 0  ⇒  V contains wave_x
σ.k_y ≠ 0  ⇒  V contains wave_y
σ.k_z ≠ 0  ⇒  V contains wave_z
σ.k_t ≠ 0  ⇒  V contains wave_t
\end{verbatim}

\paragraph{\texorpdfstring{\texttt{δV\ :\ OneDWave{[}κ{]}} (added in
v2.0)}{δV : OneDWave{[}κ{]} (added in v2.0)}}\label{ux3b4v-onedwaveux3ba-added-in-v2.0}

A difference structure isomorphic to V. Holds how V changes in one
processing step:

\begin{verbatim}
δV[j] = (
    Δθ₀ : DeltaPhase,    -- phase displacement
    ΔA  : Scalar256,      -- amplitude displacement
    Δℓ  : DeltaPhase,     -- spread displacement
)
\end{verbatim}

For free particles, δV{[}j{]}.Δθ₀ can be computed from σ:

\[\delta V[j].\Delta\theta_0 = \sigma.k_{a(j)} \cdot \delta t\]

However, in interaction events δV may change independently.

\paragraph{\texorpdfstring{\texttt{δt\ :\ DeltaT}}{δt : DeltaT}}\label{ux3b4t-deltat}

\textbf{Processing step period} (NOT time). A universal computation
counter.

\begin{itemize}
\tightlist
\item
  \(\delta t > 0\): forward operation (causal forward direction)
\item
  \(\delta t < 0\): reverse operation (guarantees reversibility)
\item
  \(\delta t = 0\): at rest
\end{itemize}

What is observed as ``time'' is the accumulation of V{[}t-axis{]}.θ₀,
not δt itself (see §6.4).

\paragraph{\texorpdfstring{\textasciitilde{}\texttt{δδt\ :\ DeltaT}
(removed in
v2.1)\textasciitilde{}}{\textasciitilde δδt : DeltaT (removed in v2.1)\textasciitilde{}}}\label{ux3b4ux3b4t-deltat-removed-in-v2.1}

The Particle structure included a \texttt{δδt} field in v2.0, but this
is removed in v2.1. Reasons:

\begin{itemize}
\tightlist
\item
  \textbf{Reversibility violation}: Executing \texttt{δt\ +=\ δδt} in
  \texttt{step(p)} does not return to the original state when δt's sign
  is reversed
\item
  \textbf{Conservation-law violation}: Single-particle δt change implies
  energy/momentum change ex nihilo, inconsistent with conservation laws
\item
  \textbf{GR free-particle principle}: Acceleration effects must arise
  from external forces (interactions); the basic operation for a free
  particle should not handle them
\end{itemize}

Acceleration effects are represented in the interaction processing
(between particles or between particle and Mesh) of §3 onward, where σ
or δt change in a conservation-consistent manner. See the separate paper
D2.

\subsubsection{§5.3 Free Particle Dynamics Close with
δV}\label{free-particle-dynamics-close-with-ux3b4v}

Although the second-order closure of physical laws (Newtonian mechanics,
GR, QM) initially suggested that δδt was needed, in v2.1:

\begin{itemize}
\tightlist
\item
  Free particles undergo only uniform geodesic motion (GR principle)
\item
  Acceleration is the result of interaction
\item
  The basic operation for a single particle is \textbf{only the update
  of phase velocity (δV) and current phase (V)}
\end{itemize}

Therefore, \textbf{4 fields (σ, V, δV, δt) form a complete
representation for free particles.}

\subsubsection{§5.4 Memory Footprint
(v2.1)}\label{memory-footprint-v2.1}

\begin{verbatim}
σ      : 6 × 256 = 1,536 bit
V      : κ × 3 × 256 ≤ 1,536 bit (κ ≤ 2)
δV     : κ × 3 × 256 ≤ 1,536 bit
δt     : 256 bit
Total  ≤ 4,864 bit = 608 byte / particle
\end{verbatim}

A 32-byte reduction from v2.0's 640 byte (removal of δδt field). Still
extremely small compared to v1.0's 224--416 byte (which lacked δV).

\begin{center}\rule{0.5\linewidth}{0.5pt}\end{center}

\subsection{§6 Operation Rules}\label{operation-rules}

\subsubsection{§6.1 Basic: 1-Step Update
(v2.1)}\label{basic-1-step-update-v2.1}

The basic operation \texttt{step(p)} consists of only one operation:

\begin{verbatim}
function step(p : Particle | null) → Particle | null:
    if p == null:
        return null
    for j in range(p.κ):
        p.V[j].θ₀ += p.σ.k_{a(j)} · p.δt
    return p
\end{verbatim}

For each non-zero position-type axis \(j\):

\[V[j].\theta_0 \mathrel{+}= \sigma.k_{a(j)} \cdot \delta t\]

That is all. \texttt{σ}, \texttt{δV}, and \texttt{δt} are not updated by
\texttt{step(p)}.

\textbf{Change from v2.0}: v2.0 included \texttt{δt\ +=\ δδt}, but this
is removed in v2.1 (see §5.2).

\subsubsection{§6.2 Type Consistency
Check}\label{type-consistency-check}

\begin{verbatim}
σ.k_axis · δt : SpatialFreq × DeltaT 
              = Linv256 × L256 
              = Scalar256 [1]                ✓

V[j].θ₀ + (σ.k · δt) : CyclePhase + Scalar256 
                     = CyclePhase [1] (mod P/2)  ✓
\end{verbatim}

\subsubsection{§6.2.1 Compliance with the Five Operational Principles
(v2.1)}\label{compliance-with-the-five-operational-principles-v2.1}

\texttt{step(p)} satisfies all five basic operational principles:

\begin{enumerate}
\def\labelenumi{\arabic{enumi}.}
\tightlist
\item
  \textbf{Null transparency}: explicitly handled at the function's start
  ✓
\item
  \textbf{In-place update}: V{[}j{]}.θ₀ updated directly, no new
  Particle generated ✓
\item
  \textbf{Type integrity}: complies with D1 §2.3 rules (see §6.2) ✓
\item
  \textbf{INPUT/OUTPUT exchangeability}: function has symmetric
  read/write structure ✓
\item
  \textbf{Time-reversal reversibility}: complete reversal by
  sign-flipping δt:
\end{enumerate}

\begin{verbatim}
First call (δt = +δt₀):  V → V + σ.k · δt₀
Sign-flip δt (δt → -δt₀)
Second call (δt = -δt₀): V → V + σ.k · δt₀ + σ.k · (-δt₀) = V  ✓
\end{verbatim}

This structurally guarantees the information-conservation requirement of
{[}W3{]}.

\subsubsection{§6.3 Polarization and Multi-Axis Phase
Difference}\label{polarization-and-multi-axis-phase-difference}

For particles with multiple non-zero axes (e.g., a polarized directional
photon):

\begin{verbatim}
σ_polarized_light = (+k_x, +k_y, 0, 0, +R_γ, 0)
V = [wave_x, wave_y]
\end{verbatim}

Polarization phase difference:

\[\Delta\varphi = V[0].\theta_0 - V[1].\theta_0 \in \text{CyclePhase}\]

{\def\LTcaptype{none} % do not increment counter
\begin{longtable}[]{@{}ll@{}}
\toprule\noalign{}
Polarization State & \(\Delta\varphi\) \\
\midrule\noalign{}
\endhead
\bottomrule\noalign{}
\endlastfoot
Linear (in-phase) & \(0\) \\
Linear (anti-phase) & \(P/4 = 180°\) \\
Circular (left) & \(+P/8 = +90°\) \\
Circular (right) & \(-P/8 = -90°\) \\
Elliptical & any integer value \\
\end{longtable}
}

All exactly representable as int256 integers (integer multiples of
\(15°\) correspond to integer multiples of \(P/48\) without error).

\subsubsection{§6.4 Treatment of Time: δt is a Processing Step, Time is
V{[}t{]}.θ₀}\label{treatment-of-time-ux3b4t-is-a-processing-step-time-is-vt.ux3b8ux2080}

\textbf{Important principle}: δt is not time but a universal processing
step. What is observed as time is the accumulation of V{[}t-axis{]}.θ₀
for each particle.

Electron example (σ.k\_t ≠ 0):

\begin{verbatim}
σ_electron = (+k_x, 0, 0, +k_t, 0, 0)
V = [wave_x, wave_t]
1-step:
  V[0].θ₀ += σ.k_x · δt   -- x-axis phase advances (momentum equivalent)
  V[1].θ₀ += σ.k_t · δt   -- t-axis phase advances ("time experienced by electron")
\end{verbatim}

Photon example (σ.k\_t = 0):

\begin{verbatim}
σ_photon = (+k_x, +k_y, 0, 0, +R_γ, 0)
V = [wave_x, wave_y]      -- no t-axis component
-- because σ.k_t = 0, no t-axis OneDWave exists in V
-- consequence: zero proper time of photon is structurally guaranteed
\end{verbatim}

\subsubsection{§6.5 Acceleration Effects Are Handled in Interaction
Processing
(v2.1)}\label{acceleration-effects-are-handled-in-interaction-processing-v2.1}

In v2.0 we stated ``δδt is determined by the background R(r) or by
interactions,'' but in v2.1:

\begin{itemize}
\tightlist
\item
  \textbf{Acceleration effects cannot be handled in the basic
  single-particle operation} (violates conservation laws and
  reversibility, see §5.2)
\item
  \textbf{All acceleration effects are realized in interaction
  processing} (where σ or δt are exchanged between particles or with the
  Mesh)
\end{itemize}

The concrete interaction rules are defined in the separate paper D2.
Gravitational time dilation, polarization rotation, and frequency shifts
are all expressed as the cumulative result of interaction events that
satisfy conservation laws.

\subsubsection{§6.6 Operations Not Required (same as
v1.0)}\label{operations-not-required-same-as-v1.0}

The following operations are not required in this specification:

\begin{enumerate}
\def\labelenumi{\arabic{enumi}.}
\tightlist
\item
  \textbf{Division}: \(R/R_1\) is converted to \(R \cdot R_1^{-1}\) as
  multiplication
\item
  \textbf{Square root}: a continuum concept; in discrete systems,
  integer values directly determine physical quantities
\item
  \textbf{Trigonometric functions}: tools of continuum sine-Gordon;
  replaceable by discrete modulation tables
\item
  \textbf{Floating-point arithmetic}: all fields are int256, no rounding
  error structurally arises
\end{enumerate}

Where cos/sin would be needed (e.g., longitudinal mode modulation),
discrete modulation tables (lookup tables) of finite resolution are
used.

\begin{center}\rule{0.5\linewidth}{0.5pt}\end{center}

\subsection{§7 Introduction of the Mesh Structure (added in
v2.0)}\label{introduction-of-the-mesh-structure-added-in-v2.0}

\subsubsection{§7.1 Separation of Longitudinal
Modes}\label{separation-of-longitudinal-modes}

In the interpretation of {[}W8{]} Table A, the Higgs
(\(\sigma_H = (0, 0, +k_z, 0, 0, 0)\)) differs from other spin-1 bosons
in having \(\sigma.R = 0\) and \(\text{spin} = 0\). Structurally it is
natural to treat the Higgs not as a ``transverse particle vibrating
position phase'' but as a \textbf{``longitudinal vibration of the mesh
itself.''}

Similarly, the graviton (gravitational waves) can be interpreted as
longitudinal-like vibration of the spacetime metric --- a vibrational
mode of the Mesh.

Treating these longitudinal modes as Particles strains the data
structure (OneDWave is essentially a type for phase rotation and cannot
directly represent axial expansion-contraction). This paper introduces a
\textbf{Mesh structure to manage longitudinal modes separately from
Particles}.

\subsubsection{§7.2 Structure of Mesh
(v2.1)}\label{structure-of-mesh-v2.1}

\begin{verbatim}
Mesh = (
    -- background scale
    R_0           : L256,             -- central length of central projection (background)
    
    -- single-axis longitudinal modes (spin 0, SO(3) symmetry breaking)
    H_x           : OneDWave,         -- x-axis longitudinal
    H_y           : OneDWave,         -- y-axis longitudinal
    H_z           : OneDWave,         -- z-axis longitudinal = Higgs
    H_t           : OneDWave,         -- t-axis longitudinal
    
    -- multi-axis longitudinal modes (spin 2, gravitational wave polarizations)
    H_xy_anti     : OneDWave,         -- xy anti-phase = gravitational wave + polarization
    H_xy_quarter  : OneDWave,         -- xy 90° phase = gravitational wave × polarization
    H_yz_anti     : OneDWave,         -- yz anti-phase
    H_xz_anti     : OneDWave,         -- xz anti-phase
    
    -- displacements of longitudinal modes
    δH_x, δH_y, δH_z, δH_t            : OneDWave,
    δH_xy_anti, δH_xy_quarter         : OneDWave,
    δH_yz_anti, δH_xz_anti            : OneDWave,
    
    -- global processing step
    δt_mesh       : DeltaT,
)
\end{verbatim}

\textbf{Changes from v2.0}:

\begin{itemize}
\tightlist
\item
  Removed \texttt{H\_xyz\_inphase} (xyz in-phase) and replaced with
  axis-specific single-axis modes \texttt{H\_x,\ H\_y,\ H\_z,\ H\_t}
  (the xyz-isotropic interpretation cannot explain the Higgs's
  generational mass hierarchy; see {[}W11{]} §8.5)
\item
  The Higgs is positioned as \texttt{H\_z} (z-axis longitudinal mode)
\item
  Removed \texttt{δδt\_mesh} (same reason as Particle: reversibility and
  conservation)
\end{itemize}

\subsubsection{§7.3 Physical Identification of Longitudinal
Modes}\label{physical-identification-of-longitudinal-modes}

{\def\LTcaptype{none} % do not increment counter
\begin{longtable}[]{@{}
  >{\raggedright\arraybackslash}p{(\linewidth - 6\tabcolsep) * \real{0.2500}}
  >{\raggedright\arraybackslash}p{(\linewidth - 6\tabcolsep) * \real{0.2500}}
  >{\raggedright\arraybackslash}p{(\linewidth - 6\tabcolsep) * \real{0.2500}}
  >{\raggedright\arraybackslash}p{(\linewidth - 6\tabcolsep) * \real{0.2500}}@{}}
\toprule\noalign{}
\begin{minipage}[b]{\linewidth}\raggedright
Mode
\end{minipage} & \begin{minipage}[b]{\linewidth}\raggedright
Vibration Pattern
\end{minipage} & \begin{minipage}[b]{\linewidth}\raggedright
Physical Identification
\end{minipage} & \begin{minipage}[b]{\linewidth}\raggedright
Spin
\end{minipage} \\
\midrule\noalign{}
\endhead
\bottomrule\noalign{}
\endlastfoot
H\_z & z-axis longitudinal (SO(3) breaking) & \textbf{Higgs} & 0 \\
H\_x, H\_y, H\_t & each-axis longitudinal & physical identification
undetermined (candidates) & 0 \\
H\_xy\_anti & x stretches while y compresses (anti-phase) &
gravitational wave + polarization & 2 \\
H\_xy\_quarter & xy 90° phase rotation & gravitational wave ×
polarization & 2 \\
H\_yz\_anti & yz anti-phase & gravitational wave (other polarization) &
2 \\
H\_xz\_anti & xz anti-phase & gravitational wave (other polarization) &
2 \\
\end{longtable}
}

\textbf{Higgs = z-axis longitudinal mode (H\_z)}: The z-axis selection
in W8 Table A's \(\sigma_H = (0, 0, +k_z, 0, 0, 0)\) is preserved as the
direction of SO(3) symmetry breaking. Third-generation fermions sharing
the z-axis have the strongest coupling and acquire the largest mass
(consistent with {[}W11{]} §8).

\textbf{Two polarizations of gravitational waves (+, ×)}: expressed as
multi-axis longitudinal modes (H\_xy\_anti, H\_xy\_quarter, etc.). SO(3)
symmetry is preserved.

\subsubsection{§7.4 Coupling between Mesh and Particle
(v2.1)}\label{coupling-between-mesh-and-particle-v2.1}

The vibrational frequency σ.k of a particle is effectively modulated by
the state of the Mesh. Mass generation and frequency shifts are
expressed as interactions through this modulation (details in the
separate paper D2).

Effective frequency (mass generation example):

\[\sigma.k_{\text{effective}}(particle, mesh) = \sigma.k(particle) - \delta k_H(\text{Mesh.H\_z VEV})\]

Third-generation fermions sharing the z-axis couple most strongly with
the H\_z VEV and acquire the largest mass (consistent with {[}W11{]}
§8).

\subsubsection{§7.5 Dynamics (v2.1: basic operation is step
only)}\label{dynamics-v2.1-basic-operation-is-step-only}

Mesh update:

\begin{verbatim}
For each longitudinal mode H_*:
    H_*.θ₀ += (intrinsic frequency) · δt_mesh
\end{verbatim}

Particle update (independent of Mesh, basic operation is
\texttt{step(p)} only):

\begin{verbatim}
For each particle p:
    step(p)   -- as in §6.1, V only is updated
\end{verbatim}

The Mesh-Particle coupling (mass generation, gravitational waves, etc.)
is not realized within single-particle \texttt{step(p)}, but in
interaction processing (defined in D2) where σ or δt are changed.

\subsubsection{§7.6 Memory Footprint
(v2.1)}\label{memory-footprint-v2.1-1}

\begin{verbatim}
R_0    : 256 bit
H_*    : 8 × 3 × 256 = 6,144 bit  -- 4 single-axis + 4 multi-axis = 8 modes
δH_*   : 8 × 3 × 256 = 6,144 bit
δt_mesh : 256 bit
Total  = 12,800 bit = 1,600 byte / entire universe
\end{verbatim}

About 1.6 KB for the entire universe. The total memory for \(N\)
particles is \(608N + 1600\) byte (a 32-byte reduction per particle from
v2.0).

\begin{center}\rule{0.5\linewidth}{0.5pt}\end{center}

\subsection{§8 Verification of W3 Structural Requirements
(v2.1)}\label{verification-of-w3-structural-requirements-v2.1}

{\def\LTcaptype{none} % do not increment counter
\begin{longtable}[]{@{}
  >{\raggedright\arraybackslash}p{(\linewidth - 4\tabcolsep) * \real{0.3333}}
  >{\raggedright\arraybackslash}p{(\linewidth - 4\tabcolsep) * \real{0.3333}}
  >{\raggedright\arraybackslash}p{(\linewidth - 4\tabcolsep) * \real{0.3333}}@{}}
\toprule\noalign{}
\begin{minipage}[b]{\linewidth}\raggedright
W3 Requirement
\end{minipage} & \begin{minipage}[b]{\linewidth}\raggedright
Satisfied
\end{minipage} & \begin{minipage}[b]{\linewidth}\raggedright
Basis
\end{minipage} \\
\midrule\noalign{}
\endhead
\bottomrule\noalign{}
\endlastfoot
Boundedness & ✓ & int256 represents all phase quantities (§1) \\
Discreteness & ✓ & all fields are integer types (§2--§5) \\
\textbf{Information conservation (time-reversal reversibility)} &
\textbf{✓ (rigorous in v2.1)} & step(p) is fully reversed by
sign-flipping δt (§6.2.1) \\
Avoidance of zero division & ✓ & division not used (§6.6) \\
Closure under integer arithmetic & ✓ & only addition/multiplication; cos
via discrete modulation table (§7.4) \\
\textbf{6-axis symmetry} & \textbf{✓ (v2.0+)} & frequency interpretation
makes all axes equivalent (§1.4, §3) \\
\textbf{Transverse-longitudinal separation} & \textbf{✓ (v2.0+)} &
Particle vs Mesh (§7) \\
\textbf{Structural conservation laws} & \textbf{✓ (rigorous in v2.1)} &
step(p) does not change σ or δt; conservation is automatic \\
Non-privileging of time & ✓ & t equal to other axes; time is V{[}t{]}.θ₀
(§6.4) \\
\textbf{All five operational principles satisfied} & \textbf{✓ (v2.1)} &
null transparency, in-place, type integrity, exchangeability,
reversibility (§6.2.1) \\
Finite representability & ✓ & ≤ 608 byte per particle, +1.6 KB for whole
universe (§5.4, §7.6) \\
\end{longtable}
}

\begin{center}\rule{0.5\linewidth}{0.5pt}\end{center}

\subsection{§9 Incompatibilities with
v1.0}\label{incompatibilities-with-v1.0}

\subsubsection{§9.1 Changes in
Interpretation}\label{changes-in-interpretation}

\begin{itemize}
\tightlist
\item
  σ.k unified from ``wave number / acceleration'' to ``vibrational
  frequency''
\item
  δt clarified from ``proper time step'' to ``processing step period''
\item
  ``Time'' is a derived quantity appearing as V{[}t-axis{]}.θ₀
\end{itemize}

\subsubsection{§9.2 Structural Extension and
Refinement}\label{structural-extension-and-refinement}

\textbf{Added in v2.0} (continued in v2.1):

\begin{itemize}
\tightlist
\item
  δV added to Particle (v1.0's Particle is included as a subset)
\item
  Mesh structure newly introduced (not present in v1.0)
\end{itemize}

\textbf{Added in v2.0 and removed in v2.1}:

\begin{itemize}
\tightlist
\item
  Particle's \texttt{δδt} field (removed for conservation/reversibility
  violations)
\item
  Mesh's \texttt{δδt\_mesh} field (same reason)
\item
  Mesh's \texttt{H\_xyz\_inphase} mode (cannot explain observed
  generational mass hierarchy; replaced with axis-specific modes such as
  \texttt{H\_z})
\end{itemize}

\subsubsection{§9.3 Compatibility with {[}W8{]} Table
A}\label{compatibility-with-w8-table-a}

The 62-particle classification of {[}W8{]} Table A remains valid in
v2.1. The Higgs is positioned as an \textbf{excitation of the Mesh's
H\_z mode} (the H\_xyz\_inphase proposed in v2.0 has been corrected in
v2.1, consistent with {[}W11{]} §8.5). The total count (61 SM states + 1
graviton = 62) is unchanged.

\begin{center}\rule{0.5\linewidth}{0.5pt}\end{center}

\subsection{§10 Conclusion}\label{conclusion}

Essence of v2.1:

\begin{itemize}
\tightlist
\item
  \textbf{xyzt are frequencies} (natural interpretation in phase space)
\item
  \textbf{6 axes are completely symmetric} (4 position-type + R + Q)
\item
  \textbf{Transverse (Particle) and longitudinal (Mesh) are separated}
\item
  \textbf{Basic operation step(p) updates only V} (satisfies all 5
  principles; acceleration effects are separated into interaction
  processing)
\item
  \textbf{Time is V{[}t{]}.θ₀} (a position phase, advancing per
  particle)
\item
  \textbf{All operations close under int256 integer arithmetic}
\end{itemize}

With this structure, the dynamics rules (D2, separate paper) can
describe:

\begin{itemize}
\tightlist
\item
  Gravitational time dilation, polarization rotation (particle-Mesh
  interaction)
\item
  Higgs mechanism (Mesh.H\_z VEV coupling to particle's σ.k\_z)
\item
  Gravitational wave emission (excitation of Mesh.H\_xy\_anti mode)
\item
  Coulomb 2-body problem (inter-particle interaction)
\end{itemize}

all within the same framework. All acceleration effects are realized
within interaction processing in a manner that respects conservation
laws and reversibility.

This paper is a design specification; these dynamical applications are
deferred to separate papers.

\begin{center}\rule{0.5\linewidth}{0.5pt}\end{center}

\subsection{References}\label{references}

\begin{itemize}
\tightlist
\item
  {[}W3{]} Kihara, N. ``Structural Requirements for a Theory of
  Everything.'' DOI: 10.5281/zenodo.19601592 (2026).
\item
  {[}W8{]} Kihara, N. ``Combinatorial Structure of the 6-Dimensional
  Hyperrectangle.'' DOI: 10.5281/zenodo.19762134 (2026).
\item
  {[}W9{]} Kihara, N. ``sine-Gordon Equation --- Foundation of
  Topological Solitons.'' DOI: 10.5281/zenodo.19650966 (2026).
\item
  {[}W10{]} Kihara, N. ``Four Modes of Shape-Invariant Waves.'' DOI:
  10.5281/zenodo.19709798 (2026).
\item
  {[}W11{]} Kihara, N. ``Interactions of Shape-Invariant Waves.'' DOI:
  10.5281/zenodo.19763463 (2026).
\item
  {[}M1{]} Kihara, N. ``Three-Layer Model of Central Projection:
  R--R₁--R₀.'' DOI: 10.5281/zenodo.19691713 (2026).
\item
  {[}TE\_Raxis{]} Kihara, N. ``Reexamination of the 6-Dimensional Sign
  Vector xyztRQ.'' DOI: 10.5281/zenodo.19904714 (2026).
\end{itemize}

\begin{center}\rule{0.5\linewidth}{0.5pt}\end{center}

\textbf{Author}: Noriaki Kihara \textbf{ORCID}:
\href{https://orcid.org/0009-0004-6753-4020}{0009-0004-6753-4020}
\textbf{License}: CC BY 4.0

\end{document}
